% ch05a —— §5.1–5.3（原文 wcf-paper.tex 2129–2544 行）

\section{相互非局域修正} \label{sec:constructing-moduli}

正如我们刚刚在最简单的非平凡情形中所见，精确的超凯勒度量 $g$ 并不等于半平坦度量 $g^{\sfm}$，因为存在来自瞬子的量子修正，而这些瞬子对应于 $d=4$ 的 BPS 态。

一般地，对每一个支撑 BPS 态的电荷 $\gamma$，我们都预期有这样一项量子修正。这些修正应以 BPS 简并度 $\degen(\gamma;u)$ 为权重。然而我们知道，当 $u$ 穿过边缘稳定墙时，$\degen(\gamma;u)$ 会发生跳变！于是似乎出现了一个困惑：$g$ 的量子修正难道不也会不连续地跳变吗？这与场论中「$g$ 应当光滑」的预期如何相容？

本节将给出大 $R$ 下精确的超凯勒度量 $g$ 的一个显式构造。我们将看到，只要 Kontsevich--Soibelman 穿墙公式成立，$g$ 确实是光滑的。这就是我们对 WCF 的物理解释。

把 $g$ 在 $R\to\infty$ 附近展开，我们便找到了上述困惑的解答：除单个 BPS 粒子的贡献之外，还有多粒子贡献。单粒子贡献中的不连续性，被多粒子部分的不连续性所补偿。关于 $\CN=1$ 语境中的相关讨论，见 \cite{GarciaEtxebarria:2007zv}。

\subsection{黎曼--希尔伯特问题的定义} \label{sec:defining-rh}

我们从第 4 节所研究的阿贝尔理论中汲取灵感，通过在 $\zeta$ 平面中求解一个黎曼--希尔伯特问题来构造该度量。% TODO 合稿改 \ref{sec:periodicnut}
我们起初在固定的 $u\in\vm$（4d 模空间）处工作，远离各边缘稳定墙。我们还在 $u$ 处取定一个二次加细 $\sigma$。

黎曼--希尔伯特问题是用映射 $\CX$ 表述的：$\CX$ 从实环面 $\widetilde\CM_u$（$\CM$ 为 3d 模空间）映入复化辛 $2r$-环面 $\tilde T_u$，后者是我们在回顾 Kontsevich--Soibelman 公式时（第 2.2 节）引入的。% TODO 合稿改 \ref{sec:ReviewKS}
给定任何这样的 $\CX$，我们把 $\tilde T_u$ 上的坐标函数 $X_\gamma$ 拉回，得到 $\widetilde\CM_u$ 上的函数 $\CX_\gamma$，其定义为 $\CX_\gamma(\theta) = X_\gamma(\CX(\theta))$。特别地，第 3.3 节所给的诸 $\CX^{\sfm}_\gamma$ 来自一个参考映射 $\CX^{\sfm}$；它就是我们下面所构造的 $\CX$ 的第零阶近似。% TODO 合稿改 \ref{sec:semiflat-twistor}

为表述这个黎曼--希尔伯特问题，我们需要确定 $\CX$ 当 $\zeta\to0,\infty$ 时的\emph{渐近行为}，以及它在 $\zeta$ 平面中的\emph{不连续性}。

我们先从渐近行为说起。引入
\begin{equation} \label{eq:def-upsilon}
\Upsilon := \CX (\CX^{\sfm})^{-1}.
\end{equation}
在本节中，我们对映射的复合采用一种非传统的记号：$(fg)(x)$ 意为 $g(f(x))$。\footnote{这种记号的一个好处，从下面的复合即可看出：$A \stackrel{f}{\to} B \stackrel{g}{\to} C$ 复合为 $A \stackrel{fg}{\to} C$。第二个好处将在第 5.5 节显现。}% TODO 合稿改 \ref{sec:differential-problem}
于是 $\Upsilon$ 是从 $\widetilde\CM_u$ 到自身的映射（更确切地说，是到其复化）。具体地，$\Upsilon$ 把
\begin{equation} \label{eq:upsilon-concrete}
e^{i \theta_i} \mapsto \CX^i(\theta) \exp\biggl[  - \pi R
\frac{Z_i}{\zeta} - \pi R  \zeta \bar Z_i \biggr].
\end{equation}
我们要求 $\Upsilon$ 当 $\zeta\to0$ 与 $\zeta\to\infty$ 时的极限存在，
\begin{equation} \label{eq:RH-asymptotics}
\lim_{\zeta \to 0} \Upsilon = \Upsilon_0, \quad \lim_{\zeta \to \infty} \Upsilon = \Upsilon_\infty,
\end{equation}
并且满足
\begin{equation} \label{eq:upsilon-reality}
\Upsilon_0 = \bar \Upsilon_\infty.
\end{equation}

接下来我们需要指定 $\CX$ 的不连续性，这里把 $\CX$ 视为 $\zeta$ 的分段解析函数。暂时假定 $u$ 不在任何边缘稳定墙上。不连续性将由与 BPS 态相伴的 Kontsevich--Soibelman 辛同胚 $\CK_\gamma:\tilde T_u\to\tilde T_u$ 给出。我们给 $\zeta$ 平面中过原点的每一条射线 $\ell$ 配以 $\clat_u$ 的一个子集
\begin{equation} \label{eq:clat-ell}
(\clat_u)_\ell := \{ \gamma: Z_\gamma(u)/\zeta \in \IR_- \text{ 对 } \zeta \in \ell \},
\end{equation}
以及相应地，对 BPS 态取的一个乘积：
\begin{equation}
S_\ell := \prod_{\gamma \in (\clat_u)_\ell} \CK_{\gamma}^{\degen(\gamma;u)}.
\end{equation}
（由于 $u$ 不在墙上，$(\clat_u)_\ell$ 至多是一维的，而诸 $\gamma\in(\clat_u)_\ell$ 对应的 $\CK_\gamma$ 彼此对易；因此无须指定这个乘积中的次序。）由于电荷格 $\clat_u$ 是可数的，除可数条射线 $\ell$ 之外，我们都有 $(\clat_u)_\ell=\emptyset$，从而 $S_\ell=1$。我们把满足 $(\clat_u)_\ell\neq\emptyset$ 的射线称为「BPS 射线」。

最典范的不连续性取法是要求
\begin{equation} \label{eq:RH-jumps}
\CX^{+} = \CX^{-} S_\ell
\end{equation}
其中 $\CX^{+}$、$\CX^{-}$ 分别是 $\zeta$ 沿顺时针、逆时针方向趋近 $\ell$ 时 $\CX$ 的极限。这是对我们在第 4 节所得结果的最直截了当的推广：在那里，我们得到在 $\zeta$ 中分段解析的映射 $\CX=(\CX_\magn,\CX_\elec)$，它有两条 BPS 射线 $\ell_\pm$；跨过 $\ell_\pm$ 时 $\CX_\elec$ 连续，而 $\CX_\magn$ 按式 (4.52) 跳变。% TODO 合稿改 \eqref{eq:X2-disc}；「第 4 节」改 \ref{sec:periodicnut}
这两条 BPS 射线对应于电荷为 $(0,\pm q)$ 的单个超多重态，其不连续性恰为 \eqref{eq:RH-jumps}，其中 $S_{\pm\ell}=\CK_{0,\pm q}$。现在我们把许多 BPS 粒子都纳入进来，办法只是要求沿许多 BPS 射线发生跳变。在这种更一般的情形下，将不存在处处连续的 $\CX_\gamma$。

至此，我们的黎曼--希尔伯特问题已经表述完毕。它的解并不唯一：确切地说，它只确定到相差一个变换 $\CX\to b\CX$，其中 $b$ 是 $\widetilde\CM_u$ 的任意微分同胚。下一节求解该问题时，我们将用一种方便的方式固定这一不确定性。

这个黎曼--希尔伯特问题或许显得有些不合常规，因为它是用 $\CX$ 与 $S_\ell$ 表述的，而它们不是线性映射，而是更一般的流形之间的映射。有顾虑的读者尽可把问题「线性化」：不去考虑 $\CX$，而考虑拉回操作 $\CX^*: C^\infty(\tilde T_u)\to C^\infty(\widetilde\CM_u)$。这样做的代价是：并非每个映射 $Q: C^\infty(\tilde T_u)\to C^\infty(\widetilde\CM_u)$ 都能写成某个映射 $\CX$ 的 $\CX^*$；因此，即使我们为线性化后的黎曼--希尔伯特问题找到了一个解 $Q$，还要额外费心验证 $Q=\CX^*$ 对某个 $\CX$ 成立。所幸黎曼--希尔伯特问题的「函子性」会来救场：$Q$ 是某个 $\CX$ 的 $\CX^*$，当且仅当它保持乘法，即 $Q(fg)=Q(f)Q(g)$。由于定义该线性问题的全部数据都与这一结构相容，其解也与之相容。

最后我们讨论本问题的一个实性（reality）性质，它对我们构造超凯勒度量至关重要。由于关系 $\degen(\gamma;u)=\degen(-\gamma;u)$，我们的不连续性条件具有一种离散对称性：给定任一解 $\CX$，我们都能得到另一个解，记作 $\bar\CX$，其定义为
\begin{equation}
\bar\CX_\gamma(\zeta) = \overline{\CX_{-\gamma}(-1/\bar\zeta)}.
\end{equation}
我们断言，事实上我们的解在这一变换下不变，即
\begin{equation} \label{eq:reality-X}
\CX = \bar \CX.
\end{equation}
为看出这一点，考虑映射 $Y=\bar\CX\CX^{-1}$。由于 $\bar\CX$ 与 $\CX$ 具有相同的不连续性，$Y$ 实际上是 $\zeta$ 的\emph{解析}函数。另一方面，由渐近条件 \eqref{eq:upsilon-reality}，当 $\zeta\to0,\infty$ 时 $Y\to1$。于是由刘维尔定理得 $Y=1$。

\subsection{KS 公式的作用} \label{sec:continuity-rh}

现在到了关键之处，这也是我们撰写本文的主要理由。假设对远离边缘稳定墙的任一固定 $u$，我们都找到一个解 $\CX$。那么它作为 $\zeta$ 的函数的行为便完全确定了：除在 BPS 射线处按 \eqref{eq:RH-jumps} 跳变外，它处处连续。但关于它作为 $u$ 的函数的行为，我们能说些什么呢？

我们的黎曼--希尔伯特问题中的 $u$ 依赖性来自两处。一处是渐近边界条件 \eqref{eq:RH-asymptotics}；这一依赖当然是连续的。另一处是不连续性处方 \eqref{eq:RH-jumps}。只要 $u$ 远离边缘稳定墙，这里的依赖也是连续的。但当我们穿过墙时，会发生什么？设 $u_w$ 为墙上的一个一般点。当 $u$ 从墙的一侧趋于 $u_w$ 时，对应于电荷 $\gamma=n\gamma_1+m\gamma_2$ 的诸 BPS 射线彼此相撞，汇合成一条射线 $\ell$。设 $A$ 为 $\CX$ 跨过这一组射线的总不连续性，
\begin{equation} \label{eq:disc-product}
A = \prod^\ccwarrow_{\substack{\gamma = n \gamma_1 + m \gamma_2 \\ m>0,\,n>0}}
\CK_\gamma^{\degen(\gamma;u)}.
\end{equation}
假定从这一侧取的极限 $\lim_{u\to u_w}\CX$ 存在，则它是如下黎曼--希尔伯特问题的解：跨过 $\ell$ 的不连续性为 $A$，而沿其他所有 BPS 射线的不连续性仍如前指定。

另一方面，我们也可以考虑从墙的另一侧取的极限 $\lim_{u\to u_w}\CX$。要使这两个极限一致，充分必要条件是：它们是\emph{同一个}黎曼--希尔伯特问题的解。于是这就要求由 \eqref{eq:disc-product} 算出的 $A$ 在墙的两侧相同。正如我们在第 2.2 节所回顾的，这正是 KS 穿墙公式的内容！\footnote{\eqref{eq:disc-product} 中的乘积取逆时针次序，而第 2.2 节中取顺时针次序；这一差别来自我们在第 \ref{sec:constructing-moduli} 节对映射复合所采用的非传统约定。}% TODO 合稿改 \ref{sec:ReviewKS}

我们于是断言：\emph{假定 BPS 简并度服从 KS 公式}，则黎曼--希尔伯特问题的解 $\CX$ 作为 $u$ 与 $\zeta$ 的函数是连续的，只有在 BPS 射线处例外。此外，跨过 BPS 射线的不连续性由一个辛同胚给出。

\subsection{求解黎曼--希尔伯特问题} \label{sec:solving-rh}

表述完黎曼--希尔伯特问题之后，我们想看到它有解，并理解其大 $R$ 行为。与第 4 节所考虑的那些简单情形不同——那里所有的 $S_\ell$ 彼此对易——这里我们无法为所求的 $\CX$ 写出显式的积分公式，只能走一条更间接的路线。% TODO 合稿改 \ref{sec:periodicnut}
我们利用这样一件事实：本问题的结构与文献 \cite{MR1213301,Cecotti:1993rm} 中所考虑的一个问题极为相似。事实上，我们的问题正是那里所考虑问题的无穷维版本。

在 \cite{Cecotti:1993rm} 中，黎曼--希尔伯特问题被改写成关于 $\Upsilon(\zeta)$ 的类比物 $\Phi$ 的一个积分方程。当 $R$ 足够大时，该方程把 $\Phi$ 描述为单位矩阵的一个小修正。因此它可以迭代求解：这既证明了当 $R$ 足够大时解的存在性，也给出第零阶近似 $\Phi=1$ 的首阶修正的显式公式。这些首阶修正直接用不连续因子表达。

这正是我们希望获得的、关于映射 $\CX$ 的那类信息。一种直接的做法，是写下 \cite{Cecotti:1993rm} 中积分方程的无穷维类比。这一做法只直接适用于线性的黎曼--希尔伯特问题，因此人们必须先过渡到上一小节末尾提到的线性问题。积分方程的解随后将给出函数空间之间的一个线性映射；如我们所述，这个线性映射将是某个映射 $\CX:\widetilde\CM_u\to\tilde T_u$ 的 $\CX^*$。

一个小问题是：倘若严格照搬 \cite{Cecotti:1993rm} 的处方，我们得到的解将满足边界条件 $\Upsilon_0=1$。而我们的构造需要选择不同的边界条件，即 \eqref{eq:upsilon-reality}，它的好处是与实性条件 $\CX_\gamma(\zeta)=\overline{\CX_{-\gamma}(-1/\bar\zeta)}$ 相容。\footnote{本节中我们把 $\CX=\CX(\zeta)$ 显式写出，即把 $\CX$ 视为随 $\zeta$ 变化的\emph{映射}，从而略去它对 $\widetilde\CM_u$ 的坐标 $\theta$ 的依赖。}幸运的是，只需对积分核稍作修改，就很容易写出一个顾及这一不同边界条件的积分方程变体。

这一策略似乎足以证明解的存在性，但它有一个重要的缺点：迭代求解的中间步骤不必具有 $\CX^*$ 的形式。若能有一种表述，使逼近格式的每一步本身都是映射 $\widetilde\CM_u\to\tilde T_u$，那将很有用。利用 $\tilde T_u$ 上的阿贝尔群结构，只要写出如下积分方程即可做到：
\begin{equation} \label{eq:X-integral-mult}
\CX_\gamma(\zeta) = \CX^{\sfm}_\gamma(\zeta) \exp \left[ \frac{1}{4 \pi
i}\sum_{\ell} \int_{\ell } \frac{d\zeta'}{\zeta'} \frac{\zeta' +
\zeta}{\zeta' - \zeta} \log \frac {\CX_\gamma(\zeta')}{(\CX S_\ell)_\gamma(\zeta')}\right]
\end{equation}
这里求和遍历各 BPS 射线 $\ell$。\eqref{eq:X-integral-mult} 的任何解都满足不连续性条件 \eqref{eq:RH-jumps}。此外，我们对积分核的选择保证该解还满足实性条件 \eqref{eq:upsilon-reality}。因此，\eqref{eq:X-integral-mult} 的解就是黎曼--希尔伯特问题的解。\footnote{注意，尽管黎曼--希尔伯特问题在 $\widetilde\CM_u$ 的微分同胚下不变，方程 \eqref{eq:X-integral-mult} 却并非如此；它的解是唯一的，而不是唯一到相差一个微分同胚。}

利用由式 (2.16) 给出的 Kontsevich--Soibelman 因子的显式形式，我们有 % TODO 合稿改 \eqref{eq:ks-symplect}
\begin{equation}
(\CX S_\ell)_\gamma = \CX_\gamma \prod_{\gamma' \in (\clat_u)_\ell} (1-
\sigma(\gamma') \CX_{\gamma'})^{\degen(\gamma';u)\langle
\gamma,\gamma' \rangle}
\end{equation}
（其中 $(\clat_u)_\ell$ 由式 \eqref{eq:clat-ell} 定义）。把它代入 \eqref{eq:X-integral-mult}，便得到关于 $\CX$ 的最终积分方程：
\begin{equation} \label{eq:X-integral-mult-explicit}
\CX_\gamma(\zeta) = \CX^{\sfm}_\gamma(\zeta) \exp \left[ -\frac{1}{4
\pi i} \sum_{\gamma'} \degen(\gamma';u) \langle \gamma,\gamma'
\rangle \int_{\ell_{\gamma'}} \frac{d\zeta'}{\zeta'} \frac{\zeta' +
\zeta}{\zeta' - \zeta} \log (1- \sigma(\gamma')
\CX_{\gamma'}(\zeta'))\right].
\end{equation}
正如我们已经提到的，方程 \eqref{eq:X-integral-mult-explicit} 是热力学贝特拟设的一种形式。见附录 E。% TODO 合稿改 \ref{app:TBA}

在附录 C 中，我们论证 \eqref{eq:X-integral-mult-explicit} 当 $R$ 足够大时有解，并描述 $u$ 远离各墙时其 $R\to\infty$ 的展开。% TODO 合稿改 \ref{app:integral-asymptotics}
第一个非平凡近似是
\begin{equation} \label{eq:X-first-approx}
\CX_\gamma(\zeta) \sim \CX^{\sfm}_\gamma(\zeta) \exp \left[ -\frac{1}{4
\pi i}  \sum_{\gamma'} \degen(\gamma';u) \langle \gamma,\gamma'
\rangle \int_{\ell_{\gamma'}} \frac{d\zeta'}{\zeta'} \frac{\zeta' +
\zeta}{\zeta' - \zeta} \log (1- \sigma(\gamma')
\CX^{\sfm}_{\gamma'}(\zeta'))\right],
\end{equation}
它本质上是我们在阿贝尔理论中找到的单瞬子修正的线性叠加。更高阶修正涉及诸 $\degen(\gamma';u)$ 的多重线性式，并带有某种 $R$ 依赖性，这种依赖性表明它们是多瞬子贡献。

我们在附录 C 中的论证，与 \cite{Cecotti:1993rm} 在有限维 $tt^*$ 语境中所给的论证密切相关。事实上，即使在有限维情形，我们的做法也使渐近分析得到简化；因此我们在附录 C 中也重新分析了那一情形。% TODO 合稿改 \ref{app:integral-asymptotics}（两处）

\subsubsection*{整体性问题}

通过求解黎曼--希尔伯特问题，我们得到了一个映射 $\CX:\widetilde\CM_u\to\tilde T_u$，它依赖于局部二次加细 $\sigma(\gamma)$ 的选取。这一选取通过不连续性 $\CK_\gamma$ 的定义影响黎曼--希尔伯特问题。然而，解 $\CX$ 对 $\sigma$ 的依赖方式很简单。回忆：对任意两个细化 $\sigma,\sigma'$，都存在某个 $c(\sigma,\sigma')\in\clat_u^*/2\clat_u^*$，使得 $\sigma(\gamma)\sigma'(\gamma)=(-1)^{\gamma\cdot c(\sigma,\sigma')}$。给定细化 $\sigma$ 下 \eqref{eq:X-integral-mult} 的一个解 $\CX^{[\sigma]}$，就有细化 $\sigma'$ 下相应的解 $\CX^{[\sigma']}$：
\begin{equation}
\CX_\gamma^{[\sigma']}(u,\theta;\zeta) = (-1)^{\gamma \cdot c(\sigma,\sigma')} \CX_\gamma^{[\sigma]}(u,\theta + c \pi;\zeta).
\end{equation}
由此可见，若利用细化把 $\widetilde\CM_u\simeq\CM_u$ 与 $\tilde T_u\simeq T_u$ 等同起来，我们就得到与细化选取\emph{无关}的 $\CX:\CM_u\to T_u$。
